\scanpage{79}
Je désigne par $\widehat M^*_{\underline G,\underline G'}$, $\overset{\deleted{\unsure{\wedge}}}{M}{}^*_{\underline G,\underline G'}$ les germes de multiplicités obtenus en \deleted{enlevant (``l'infini'') $\widehat M_{\underline G}$ resp. $M_{\underline G}$, $\unsure{\Gamma^!_{\underline G,\underline G'}}$ (les traces de la \unsure{réunion} des diviseurs de $\widehat M_{\underline G'}$ \unsure{qui contiennent} $(\widehat M_{\underline G},\Gamma^!_{\underline G,\underline G'})$)}
\inserted{\deleted{prenant les traces de $M$} enlevant les traces des diviseurs à l'infini de $\widehat M_{\underline G'}$ \deleted{\unsure{non seuls}} qui contiennent $M_{\underline G}$ -- dans le cas de $M^*_{\underline G,\underline G'}$, on prend la \uline{trace} de $M_{\underline G'}$ sur $\widehat M_{\underline G,\underline G'}$.}
Ils sont donc des filtres d'ouverts de la multiplicité $M_{\underline G'}$, et les groupoïdes \unsure{fondamentals|fondamentaux} \deleted{$\unsure{\Pi}$}
\[
(100)\qquad\unsure{\Pi M^*_{\underline G,\underline G'}|\Pi\widehat M^*_{\underline G,\underline G'}},\qquad\Pi M^*_{\underline G,\underline G'}
\]
sont définis comme des Catégories $\underset{\longrightarrow}{\lim}$ des catégories de revêtements \unsure{universels}, ou comme des catégories des revêtts \unsure{du connexe} desdits germes de multiplicités (comparer [5], n°5, p. $768'$).

On a alors des diagrammes de germes de multiplicités analytiques
\diagram{\begin{tikzcd}[ampersand replacement=\&, column sep=large,row sep=large]
\& M^*_{\underline G,\underline G'}\arrow[r,dash]\arrow[d] \& M_{\underline G'}\deleted{\unsure{M}}\\
(M_{\underline G}\deleted{\unsure{,\Gamma^!}})\arrow[r,"{\text{(équiv. d'homotopie !)}}"'] \& M_{\underline G,\underline G'} \&
\end{tikzcd}\quad(101)}
donnant lieu à des morphismes de groupoïdes
\margin{78}
